\section{Results} \label{sec:results}
Notizen, um das Kapitel zu schreiben:
\\
Aufgrund dass unsere Daten recht schlecht waren und wir die erst in der letzten Woche bereinigt haben, nachdem wir viele verschiedene Methoden ausprobiert haben, hatten wir nur noch begrenzt Zeit zu experimentieren.
Deswegen ist der Output von den Modellen allgemein immer noch nicht besonders gut, was aber auch an der Grösse liegen kann.

\begin{itemize}
    \item Grundidee: Bei allen wichtigen Experimenten nicht nur Trainings- und Validation-Loss betrachten,
    sondern zusätzlich Evaluationsmetriken berichten.
    \item Ziel: Bessere Vergleichbarkeit und realistischere Einschätzung der Modellqualität.
    \item Tokenizer vocab size vergleichen 
\end{itemize}


\subsection*{Baselines}
\begin{itemize}
    \item Bigram-Modell als einfache Baseline.
    \item Transformer-Baseline mit sinusoidalem Positional Embedding.
\end{itemize}


\subsection*{Vergleich Positional Embeddings}
\begin{itemize}
    \item Vergleich zweier Transformer-Varianten:
    \begin{itemize}
        \item Transformer mit sinusoidalem Positional Embedding
        \item Transformer mit RoPE (Rotary Positional Embedding)
    \end{itemize}
    \item Ergebnis: RoPE liefert bessere Ergebnisse.
    \item Konsequenz: RoPE wird in allen weiteren Experimenten verwendet.
\end{itemize}


\subsection*{Gewichtsinitialisierung}
\begin{itemize}
    \item Vergleich von Modellen:
    \begin{itemize}
        \item mit Gewichtsinitialisierung
        \item ohne Gewichtsinitialisierung
    \end{itemize}
    \item Durchführung für Transformer mit Sinusoidal Embedding und mit RoPE.
    \item Ergebnis: Gewichtsinitialisierung verbessert die Resultate.
    \item Konsequenz: Gewichtsinitialisierung wird überall verwendet.
\end{itemize}


\subsection*{Feedforward-Dimension}
\begin{itemize}
    \item Untersuchung des Einflusses der Feedforward-Hidden-Dimension.
    \item Vergleich folgender Varianten:
    \begin{itemize}
        \item $4 \times$ Embedding-Dimension (typischer Standard)
        \item $2 \times$ Embedding-Dimension
        \item $1 \times$ Embedding-Dimension
    \end{itemize}
\end{itemize}


\subsection*{Hyperparameter-Sweep (Architektur)}
\begin{itemize}
    \item Sweep über harte Modellparameter:
    \begin{itemize}
        \item Embedding-Dimension
        \item Anzahl der Blöcke (Layer)
        \item Anzahl der Attention-Köpfe
        \item Sequenzlänge
        \item Feedforward-Hidden-Dimension
    \end{itemize}
    \item Aus Zeitgründen keine Anpassung von Trainingsparametern.
    \item Nicht variiert:
    \begin{itemize}
        \item Learning Rate
        \item Batch Size
        \item Dropout
    \end{itemize}
\end{itemize}


\subsection*{GRU-Experimente}
\begin{itemize}
    \item Zusätzlich Training eines GRU-Modells.
    \item Ergebnis: GRU liefert insgesamt schlechtere Resultate als die Transformer-Modelle.
    \item Wir können zusätzlich einen Sweep für die GRU durchführen:
    \begin{itemize}
        \item Input-Dimension
        \item Hidden-Dimension
        \item Anzahl der Layer
    \end{itemize}
\end{itemize}
